% !TeX root = ./main.tex

\chapter{АРХИТЕКТУРА CDQuad И СЕРТИФИКАЦИЯ ОШИБКИ}

\section{Информационно-функциональная схема конвейера и модульная архитектура комплекса CDQuad}

Разработанный автором программно-алгоритмический комплекс \texttt{CDQuad} (\emph{Continuum-Consistent Convection-Diffusion Quadrature}) реализует сквозной конвейер обработки пространственно-временных данных от грубых бинированных кадров обучения до выдачи сертифицированных морфометрических дескрипторов на сетках произвольного разрешения. Векторная архитектурно-информационная схема взаимодействия шести функциональных подсистем комплекса представлена на рисунке~\ref{fig:tikz_pipeline}.

\begin{figure}[htbp]
\centering
\resizebox{0.98\textwidth}{!}{%
\begin{tikzpicture}[
    node distance=0.8cm and 0.9cm,
    box/.style={
        rectangle, rounded corners=5pt, draw=#1, line width=1.1pt,
        fill=#1!8, text width=4.5cm, minimum height=2.3cm,
        align=left, font=\footnotesize\sffamily, inner sep=6pt
    },
    arr/.style={-{Stealth[length=2.8mm, width=2.0mm]}, line width=1.0pt, draw=bsuBlue}
]
    \node[box=bsuBlue] (in) {
        \textbf{1. Входные данные}\\
        $\bullet$ Грубые кадры биннинга\\
        $\bullet$ Кадры полного разрешения\\
        $\bullet$ Сенсорные облака точек
    };
    \node[box=bsuGreen, right=of in] (mesh) {
        \textbf{2. Модуль \texttt{meshes.py}}\\
        $\bullet$ Триангуляция Делоне\\
        $\bullet$ Ячейки Дональда $|D_i|=m_i$\\
        $\bullet$ Матрица подъёма $P_G$ (\texttt{CSR})
    };
    \node[box=bsuAmber, right=of mesh] (op) {
        \textbf{3. Ядро \texttt{operators.py}}\\
        $\bullet$ $P_1$-жёсткость $W_{ij}(K)$\\
        $\bullet$ Upwind-потоки $F_{ij}$\\
        $\bullet$ Решатель \texttt{SciPy} LU / IMEX
    };

    \node[box=bsuAccent, below=1.0cm of in] (fit) {
        \textbf{4. Идентификация $\hat{\theta}$}\\
        $\bullet$ Оценка 6 физических параметров среды только на грубой сетке обучения
    };
    \node[box=bsuAmber, right=of fit] (cert) {
        \textbf{5. Модуль \texttt{certificate.py}}\\
        $\bullet$ Словарь 15 мод Неймана\\
        $\bullet$ Двойственный дефект $\widehat{\delta}_{-1}(G)$\\
        $\bullet$ Конформный квантиль $\widehat{q}_{1-\alpha}$
    };
    \node[box=bsuGreen, right=of cert] (desc) {
        \textbf{6. Модуль \texttt{descriptors.py}}\\
        $\bullet$ Площадь $\calM_0$, барицентр $\calM_1$\\
        $\bullet$ Интервал $\pm L_\Phi \widehat{U}_{1-\alpha}(G)$\\
        $\bullet$ Диагностика сбоев SLIC
    };

    \draw[arr] (in) -- (mesh);
    \draw[arr] (mesh) -- (op);
    \draw[arr] (op.south) -- ++(0,-0.45) -| (fit.north);
    \draw[arr] (fit) -- (cert);
    \draw[arr] (mesh) -- (cert);
    \draw[arr] (op) -- (desc);
    \draw[arr] (cert) -- (desc);
\end{tikzpicture}%
}
\caption{Архитектурно-информационная схема многомасштабного программного комплекса \texttt{CDQuad}}
\label{fig:tikz_pipeline}
\end{figure}

Основной принцип проектирования конвейера состоит в полном отделении параметров физического закона среды, определённых инвариантно на непрерывной области $\Omega=(0,1)^2$, от метрических атрибутов конкретной расчётной сетки. Программная реализация выполнена на языке Python~3.11 в виде четырёх слабосвязанных модулей пакета \texttt{src/}:
\begin{enumerate}
    \item \texttt{src/meshes.py}~-- конструирование регулярных решёток биннинга, возмущённых и центроидально-сглаженных по Ллойду триангуляций Делоне (\texttt{scipy.spatial.Delaunay}~[\ref{itm:scipy}]), вырожденных облаков датчиков и графов смежности суперпикселей SLIC (\texttt{skimage.segmentation.slic}~[\ref{itm:skimage}]), расчёт масс барицентрических ячеек Дональда $m_i$, показателя регулярности формы элементов $\rho(\calT_h)$ и кэширование барицентрических матриц подъёма $P_G$ в сжатом строчном формате \texttt{CSR} на единую проверочную квадратурную решётку $160\times 160$ ($25\,600$ узлов средних точек);
    \item \texttt{src/operators.py}~-- полностью векторизованная сборка матриц анизотропной жёсткости и дуальных конвективных потоков с индексацией неориентированных рёбер 64-битными ключами через \texttt{np.searchsorted} без циклов интерпретатора Python, блочное кэширование матрицы пространственного гауссова ядра порциями по $400$ строк для ограничения пикового потребления оперативной памяти, а также интегрирование по времени безусловно устойчивой схемой Кранка--Николсон второго порядка с однократной предфакторизацией \texttt{scipy.sparse.linalg.factorized} (в линейном режиме) либо полунеявной схемой IMEX-BDF1 с неявным учётом диффузии и реакции (в нелинейном адаптивном режиме);
    \item \texttt{src/certificate.py}~-- аналитическое зондирование разреженного оператора на словаре из $15$ собственных косинусных мод задачи Неймана в эллиптической двойственной норме $H_G^{-1}$ и непараметрическая split-conformal калибровка доверительного множителя;
    \item \texttt{src/descriptors.py}~-- вычисление глобально липшицевых морфометрических функционалов структуры, аналитический перенос границ ошибки поля на дескрипторы и формирование диагностических флагов потери симплициальной согласованности.
\end{enumerate}

Межмодульный программный интерфейс (API) организован через пару неизменяемых структур данных (\texttt{@dataclass}), исключающих утечку сеточных артефактов в вектор обучаемых параметров. Контейнер \texttt{PhysicalLawParams} хранит исключительно континуальные характеристики среды на $\Omega$: собственные коэффициенты диффузии $(\kappa_1,\kappa_2)$, функцию ориентации волокон $\theta(x)$, потоковую функцию скорости $\psi(x)$ (либо векторное поле $v(x)$), коэффициент реакции $\lambda_r$ и параметры нелокального вентиля $(\alpha_{\max}, r_0, \gamma)$. Независимый контейнер \texttt{MeshData} инкапсулирует координаты узлов $X\in\Real^{N\times 2}$, массив треугольников $\calT_h$, вектор масс ячеек Дональда $m\in\Real^N$, диаметр шага $h_G$, показатель вырождения формы $\rho(\calT_h)$ и предвычисленный \texttt{CSR}-оператор подъёма $P_G$. При переносе модели между разрешениями экземпляр \texttt{PhysicalLawParams} передаётся неизменным в процедуры сборки матриц и расчёта сертификата на любом новом объекте \texttt{MeshData}.

\section{Математическое обеспечение континуально-согласованного оператора CDQuad}

В основу вычислительного ядра положена начально-краевая задача конвекции-диффузии-реакции в дивергентной форме в цилиндре $\Omega\times(0,T]$, $\Omega=(0,1)^2$, с естественным условием непроницаемости границы $\partial\Omega$ для полного потока вещества:
\begin{equation}\label{eq:oit_pde}
\partial_t u = \divop\!\Bigl(\bigl(1-\alpha_{r_0}[u]\bigr)K(x)\nabla u\Bigr) - \divop\!\Bigl(\alpha_{r_0}[u]\,v(x)\,u\Bigr) - \lambda_r u,
\end{equation}
где $K(x)=R(\theta(x))^\top\diag(\kappa_1,\kappa_2)R(\theta(x))\in\Real^{2\times 2}$~-- симметричный равномерно положительно определённый тензор анизотропной диффузии с собственными значениями $0<\kappa_{\min}\le\kappa_{\max}<\infty$ и углом ориентации волокон среды $\theta(x)$; $v(x)=\nabla^\perp\psi(x)$~-- гладкое бездивергентное поле скорости ($\divop v\equiv 0$, $v\cdot\nu|_{\partial\Omega}=0$); $\lambda_r\ge 0$~-- коэффициент линейной кинетики затухания (фотовыцветания флуорофора); $\alpha_{r_0}[u](x)\in[0,\alpha_{\max}]$ ($\alpha_{\max}<1$)~-- нелокальный вентиль гомогенности в физическом окне радиуса $r_0>0$. Интенсивность анизотропии и относительный вклад сноса количественно определяются числом обусловленности тензора $\chi(K)$ и анизотропным числом Пекле $\Pe$:
\begin{equation}\label{eq:oit_peclet}
\chi(K) := \frac{\kappa_{\max}}{\kappa_{\min}} \ge 1, \qquad \Pe := \frac{\|v\|_{L^\infty(\Omega)}\,\diam(\Omega)}{\kappa_{\min}}.
\end{equation}

На конформной триангуляции $\calT_h=(V,E)$ с пространством кусочно-линейных базисных функций Куранта $\{\phi_i\}_{i=1}^N$ соединение барицентров треугольников $x_T$ с серединами сторон $x_e$ образует вокруг каждой вершины $i$ замкнутую барицентрическую дуальную ячейку Дональда $D_i$ площади $m_i=|D_i|=\frac{1}{3}\sum_{T\ni i}|T|=\int_\Omega\phi_i\,dx$~[\ref{itm:baba}, с.~6]. Рёберные проводимости диффузии $W_{ij}(K)$ и антисимметричные конвективные потоки $F_{ij}$ через двухзвенную ломаную границу $\Gamma_{ij}^*:=\partial D_i\cap\partial D_j$ вычисляются по 3-точечной квадратуре Гаусса--Лежандра 5-го порядка точности:
\begin{equation}\label{eq:oit_weights}
\begin{aligned}
W_{ij}(K) &:= -\int_\Omega \bigl(\nabla\phi_i(x)\bigr)^\top K(x)\nabla\phi_j(x)\,dx,\\
F_{ij} &:= \int_{\Gamma_{ij}^*} v(s)\cdot\nu_{ij}(s)\,ds = -F_{ji}.
\end{aligned}
\end{equation}
При неостроугольности триангуляции в метрике $K^{-1}$ проводимости неотрицательны ($W_{ij}(K)\ge 0$), что гарантирует дискретный принцип максимума и равномерную коэрцитивность диффузионной квадратичной формы.

Для устранения вырождения $\Theta(h_G^2)$ одношагового комбинаторного вентиля дипломной модели AdaCDE~[\ref{itm:a3}, с.~32] дискретный переключатель гомогенности строится как взвешенная площадями ячеек $m_j$ квадратура гауссова ядра $K_{r_0}(x_i,x_j)=\exp(-\|x_i-x_j\|_2^2/r_0^2)$ фиксированного физического радиуса $r_0>0$~[\ref{itm:calvello}]:
\begin{equation}\label{eq:oit_gate}
\begin{aligned}
h_{G,i}(u_G) &:= \frac{\sum_{j=1}^N m_j K_{r_0}(x_i,x_j)\exp\!\bigl(-\gamma(u_{G,i}-u_{G,j})^2\bigr)}{\sum_{j=1}^N m_j K_{r_0}(x_i,x_j)},\\
\alpha_{G,i}(u_G) &:= \alpha_{\max}\bigl(1-h_{G,i}(u_G)\bigr), \qquad \alpha_{G,ij}:=\frac{\alpha_{G,i}+\alpha_{G,j}}{2}.
\end{aligned}
\end{equation}
При шаге сетки $h_G\le \frac{1}{4}\kappa_{r_0}r_0$ оператор~\eqref{eq:oit_gate} аппроксимирует непрерывный функционал $\alpha_{r_0}[u]$ с первым порядком точности $O(h_G)$ и удовлетворяет глобальному условию Липшица в $\ell^2(V,m)$ с константой, не зависящей от числа узлов $N$. Подстановка~\eqref{eq:oit_weights} и~\eqref{eq:oit_gate} в баланс дуального объёма с монотонным противопоточным расщеплением $F_{ij}^\pm:=\max(\pm F_{ij},0)$~[\ref{itm:eymard}, с.~765] порождает разреженную систему ОДУ комплекса \texttt{CDQuad}:
\begin{multline}\label{eq:oit_ode}
m_i\frac{d u_{G,i}}{dt} = \sum_{j\in\calN_i}\bigl(1-\alpha_{G,ij}(u_G)\bigr)W_{ij}(K)(u_{G,j}-u_{G,i}) - {}\\
- \sum_{j\in\calN_i}\alpha_{G,ij}(u_G)\bigl(F_{ij}^+ u_{G,i} - F_{ij}^- u_{G,j}\bigr) - \lambda_r m_i u_{G,i}.
\end{multline}
По формуле Остроградского--Гаусса для поля $v=\nabla^\perp\psi$ суммарный поток через замкнутый контур ячейки Дональда тождественно равен нулю ($\sum_{j\in\calN_i}F_{ij}=0$). Благодаря этому кососимметричная часть переноса в энергетическом балансе обнуляется, а противопоточная стабилизация вносит дополнительную диссипацию $-\frac{1}{2}\sum_{\{i,j\}\in E}|F_{ij}|(\xi_j-\xi_i)^2\le 0$. В результате расхождение траекторий между двумя допустимыми триангуляциями $G_1,G_2$ оценивается величиной $O\bigl((\chi(K)+\Pe)(h_{G_1}+h_{G_2})\bigr)$ с равномерным по $t\in[0,\infty)$ показателем затухания $\mu=\frac{\pi^2}{4}\kappa_{\min}+2\lambda_r>0$, то есть число Пекле $\Pe$ полностью выведено из показателя экспоненты Гронуолла~[\ref{itm:a1}].

При численном интегрировании системы~\eqref{eq:oit_ode} с шагом $\Delta t = T/K_t$ в линейном режиме ($\alpha_{\max}=0$) применяется безусловно $A$-устойчивая схема Кранка--Николсон $(M_G - \frac{\Delta t}{2}L_G)u_G^{n+1} = (M_G + \frac{\Delta t}{2}L_G)u_G^n$, тождественно сохраняющая полную дискретную массу $\sum_{i=1}^N m_i u_{G,i}^n$ при $\lambda_r=0$, а в адаптивном режиме~-- полунеявная схема IMEX-BDF1, в которой жёсткий диффузионно-реакционный блок обращается неявно при явном вычислении нелокального вентиля на текущем слое. Оценка вектора параметров среды $\hat{\theta}=(\hat{\kappa}_1,\hat{\kappa}_2,\hat{\theta}_0,\hat{u}_{\mathrm{vel}},\hat{\lambda}_r,\hat{\alpha}_{\max})$ на грубой обучающей решётке биннинга сводится к минимизации взвешенного квадратичного функционала невязки траекторий $\sum_k \|u_G(t_{k+1};\theta)-\tilde{u}_G^{\mathrm{obs}}(t_{k+1})\|_{\ell^2(m)}^2$ квазиньютоновским алгоритмом с барьерными ограничениями \texttt{L-BFGS-B} (\texttt{scipy.optimize}), требующим лишь нескольких десятков прямых прогонов на малой сетке ($N\approx 300\text{--}900$).

\section{Эллиптический двойственный дефект и раздельная конформная сертификация}

На нерегулярных триангуляциях площади и ориентации граней ячеек Дональда лишены центральной симметрии, поэтому наивная поточечная $\ell^2(m)$-норма невязки противопоточного оператора имеет порядок $O(1)$ и не отражает реальное уменьшение ошибки при сгущении сетки. Чтобы согласовать индикатор качества сетки с энергетической оценкой переноса, функциональная невязка $d_G(\psi)$ дискретного оператора на словаре $\calT_{15}=\{\cos(\pi p x_1)\cos(\pi q x_2): 0\le p,q\le 3,\;(p,q)\neq(0,0)\}$ гладких мод Неймана (с аналитическим вычетом граничного конормального интеграла) измеряется в двойственной соболевской норме $H_G^{-1}$ через симметричную положительно определённую матрицу Рисса $E_G := M_G + \kappa_{\min}^{-1}S_G(K)$:
\begin{equation}\label{eq:oit_dual}
\begin{aligned}
\widehat{\delta}_{-1}(G;\psi) &:= \frac{\bigl\langle d_G(\psi),\, E_G^{-1}d_G(\psi)\bigr\rangle^{1/2}}{\|\psi\|_{H^2(\Omega)}},\\
\widehat{\delta}_{-1}(G) &:= \left(\frac{1}{15}\sum_{\psi\in\calT_{15}}\widehat{\delta}_{-1}(G;\psi)^2\right)^{1/2}.
\end{aligned}
\end{equation}
По калибровочной выборке из $B$ независимых сеток $\{G_b\}_{b=1}^B$ рассчитываются нормированные невязки $s_b = e_b / \widehat{B}(G_b;T)$, где $\widehat{B}(G;T)=\widehat{\delta}_{-1}(G)\sqrt{(1-e^{-\widehat{\mu}T})/\widehat{\mu}}$, и эмпирический квантиль $\widehat{q}_{1-\alpha}$ ранга $\lceil(B+1)(1-\alpha)\rceil$~[\ref{itm:vovk}, с.~42;~\ref{itm:angelopoulos}, с.~505], задающий верхнюю границу ошибки на произвольной тестовой сетке без обращения к эталонному решателю:
\begin{equation}\label{eq:oit_cp}
\begin{aligned}
\widehat{U}_{1-\alpha}(G_{\mathrm{test}};T) &:= \widehat{q}_{1-\alpha}\,\widehat{B}(G_{\mathrm{test}};T),\\
\mathbb{P}\bigl(e_{\mathrm{test}} \le \widehat{U}_{1-\alpha}\bigr) &\ge 1 - \alpha - d_{\mathrm{TV}}\!\bigl(s_\#\calP_{\mathrm{cal}},\,s_\#\calP_{\mathrm{test}}\bigr),
\end{aligned}
\end{equation}
где поправка на сдвиг распределения сеток контролируется расстоянием полной вариации $d_{\mathrm{TV}}$~[\ref{itm:millard}].

Поскольку классическая бинарная пороговая сегментация $A_\theta^{\mathrm{hard}}(u)=\int_\Omega\mathbf{1}_{\{u>\theta\}}\,dx$ разрывна в $L^2(\Omega)$ (малое возмущение плато интенсивности вблизи порога $\theta$ вызывает единичный скачок площади), для метрологически корректной морфометрии вводится логистическое сглаживание $\sigma_\eta(s)=(1+e^{-(s-\theta)/\eta})^{-1}$ с фиксированной физической шириной переходного слоя $\eta>0$, обеспечивающее точные аналитические константы Липшица для площади и барицентра структуры:
\begin{equation}\label{eq:oit_morph}
\begin{aligned}
\calM_0(u) &:= \int_\Omega \sigma_\eta\bigl(u(x)\bigr)\,dx, \qquad L_{\calM_0} = \frac{1}{4\eta},\\
\calM_1(u) &:= \frac{1}{\calM_0(u)}\int_\Omega x\,\sigma_\eta\bigl(u(x)\bigr)\,dx, \qquad L_{\calM_1} = \frac{\sqrt{2}}{4\eta m_0}.
\end{aligned}
\end{equation}

В таблице~\ref{tab:complexity} приведены теоретические оценки асимптотической сложности и измеренное время выполнения всех стадий конвейера.

\begin{table}[htbp]
\caption{Вычислительная сложность и время работы модулей комплекса \texttt{CDQuad} (CPU Apple Silicon, один поток)}\label{tab:complexity}
\footnotesize
\setlength{\tabcolsep}{4.5pt}
\renewcommand{\arraystretch}{1.18}
\begin{tabularx}{\textwidth}{@{} Y >{\raggedright\arraybackslash}p{2.5cm} >{\centering\arraybackslash}p{2.8cm} >{\raggedleft\arraybackslash}p{2.2cm} >{\raggedleft\arraybackslash}p{2.3cm} @{}}
\toprule
\textbf{Стадия информационного конвейера} & \textbf{Модуль} & \textbf{Сложность} & \textbf{$N=361$ ($18\times 18$)} & \textbf{$N=3025$ ($54\times 54$)} \\
\midrule
Триангуляция Делоне и массы ячеек $m_i$ & \texttt{meshes.py} & $O(N\log N)$ & 0,4~мс & 2,9~мс \\
Векторизованная сборка $W_{ij}(K)$ и $F_{ij}$ & \texttt{operators.py} & $O(|E|)=O(N)$ & 0,3~мс & 1,8~мс \\
Интегрирование Кранка--Николсон (\texttt{SciPy} LU) & \texttt{operators.py} & $O(N^{1{,}2}+K_t N)$ & 1,9~мс & 14,6~мс \\
Двойственный дефект $\widehat{\delta}_{-1}(G)$ по 15 модам & \texttt{certificate.py} & $O(N^{1{,}2}+15N)$ & 0,6~мс & 4,2~мс \\
Дескрипторы $\calM_0, \calM_1, \mathcal{V}$ и интервалы & \texttt{descriptors.py} & $O(M_{\mathrm{eval}})$ & 0,5~мс & 0,8~мс \\
\bottomrule
\end{tabularx}
\end{table}

Как следует из таблицы~\ref{tab:complexity}, благодаря кэшированию разреженной LU-факторизации для всех 15 пробных мод полный цикл расчёта прогноза и его статистической сертификации на сетке из $3025$ узлов занимает суммарно $24{,}3$~мс (из которых на сам модуль сертификации приходится лишь $4{,}2$~мс), что полностью удовлетворяет требованиям интерактивной биомедицинской диагностики в реальном масштабе времени.

\newpage
